\chapter{El problema de identificación}
La regla de cadena se evaluó distinguiendo $\beta_{within}$, pendiente medida dentro de una posición fija con efectos fijos por posición, y $\beta_{between}$, pendiente de las medias entre posiciones. La curva \emph{between} se ajustó con especificaciones lineal y pol2.

Los valores completos de pendientes, brechas, residuos, pruebas F y validación leave-one-out están en los CSV de `step5/` y `step6\_v2/step5/`. Este capítulo usa ambas campañas explícitamente: `step5/` corresponde a la campaña constante y `step6\_v2/step5/` a la campaña óptica corregida.

Las comparaciones entre especificaciones se reportan con sus chi-cuadrados, grados de libertad, valores p, RMSE y PRESS. La redacción mantiene la formulación ``compatible con'' cuando el archivo fuente no permite una identificación causal.
